\section{The extended alternating {\v C}ech complex}
\label{section-alternating-cech}

\noindent
Let $R$ be a ring. Let $f_1, \ldots, f_r \in R$. The extended alternating
{\v C}ech complex of $R$ is the cochain complex
$$
R \to \bigoplus\nolimits_{i_0} R_{f_{i_0}} \to
\bigoplus\nolimits_{i_0 < i_1} R_{f_{i_0}f_{i_1}} \to
\ldots \to R_{f_1\ldots f_r}
$$
where $R$ is in degree $0$, the term $\bigoplus_{i_0} R_{f_{i_0}}$
is in degre $1$, and so on. The maps are defined as follows
\begin{enumerate}
\item The map $R \to \bigoplus\nolimits_{i_0} R_{f_{i_0}}$
is given by the canonical maps $R \to R_{f_{i_0}}$.
\item Given $1 \leq i_0 < \ldots < i_{p + 1} \leq r$ and
$0 \leq j \leq p + 1$ we have the canonical localization map
$$
R_{f_{i_0} \ldots \hat f_{i_j} \ldots f_{i_{p + 1}}} \to
R_{f_{i_0} \ldots f_{i_{p + 1}}}
$$
\item The differentials use the canonical maps of (2) with sign $(-1)^j$.
\end{enumerate}
If $M$ is any $R$-module, the extended alternating {\v C}ech complex
of $M$ is the similarly constructed cochain complex
$$
M \to \bigoplus\nolimits_{i_0} M_{f_{i_0}} \to
\bigoplus\nolimits_{i_0 < i_1} M_{f_{i_0}f_{i_1}} \to
\ldots \to M_{f_1\ldots f_r}
$$
where $M$ is in degree $0$ as before.

\begin{lemma}
\label{lemma-extended-alternating-is-complex}
The extended alternating {\v C}ech complexes defined above
are complexes of $R$-modules.
\end{lemma}

\begin{proof}
Omitted.
\end{proof}

\begin{lemma}
\label{lemma-extended-alternating-form-module}
Let $R$ be a ring. Let $f_1, \ldots, f_r \in R$. Let $M$ be an $R$-module.
The extended alternating {\v C}ech complex of $M$ is the tensor product
over $R$ of $M$ with the extended alternating {\v C}ech complex of $R$.
\end{lemma}

\begin{proof}
Omitted.
\end{proof}

\begin{lemma}
\label{lemma-extended-alternating-base-change}
Let $R$ be a ring. Let $f_1, \ldots, f_r \in R$. Let $M$ be an $R$-module.
Let $R \to S$ be a ring map, denote $g_1, \ldots, g_r \in S$ the images
of $f_1, \ldots, f_r$, and set $N = M \otimes_R S$.
The extended alternating {\v C}ech complex constructed using
$S$, $g_1, \ldots, g_r$, and $N$ is the tensor product of the
extended alternating {\v C}ech complex of $M$ with $S$ over $R$.
\end{lemma}

\begin{proof}
Omitted.
\end{proof}

\begin{lemma}
\label{lemma-extended-alternating-homotopy-zero}
Let $R$ be a ring. Let $f_1, \ldots, f_r \in R$. Let $M$ be an $R$-module.
If there exists an $i \in \{1, \ldots, r\}$ such that $f_i$ is a unit, then
the extended alternating {\v C}ech
complex of $M$ is homotopy equivalent to $0$.
\end{lemma}

\begin{proof}
We will use the following notation: a cochain $x$ of degree $p + 1$
in the extended alternating {\v C}ech complex of $M$ is
$x = (x_{i_0 \ldots i_p})$ where $x_{i_0 \ldots i_p}$ is in
$M_{f_{i_0} \ldots f_{i_p}}$. With this notation we have
$$
d(x)_{i_0 \ldots i_{p + 1}} =
\sum\nolimits_j (-1)^j x_{i_0 \ldots \hat i_j \ldots i_{p + 1}}
$$
As homotopy we use the maps
$$
h : \text{cochains of degree }p + 2 \to \text{cochains of degree }p + 1
$$
given by the rule
$$
h(x)_{i_0 \ldots i_p} = 0 \text{ if } i \in \{i_0, \ldots, i_p\}
\text{ and }
h(x)_{i_0 \ldots i_p} = 
(-1)^j x_{i_0 \ldots i_j i i_{j + 1} \ldots i_p} \text{ if not}
$$
Here $j$ is the unique index such that $i_j < i < i_{j + 1}$ in the
second case; also, since $f_i$ is a unit we have the equality
$$
M_{f_{i_0} \ldots f_{i_p}} =
M_{f_{i_0} \ldots f_{i_j} f_i f_{i_{j + 1}} \ldots f_{i_p}}
$$
which we can use to make sense of thinking of
$(-1)^j x_{i_0 \ldots i_j i i_{j + 1} \ldots i_p}$
as an element of $M_{f_{i_0} \ldots f_{i_p}}$.
We will show by a computation that $d h + h d$ equals
the negative of the identity map which finishes the proof.
To do this fix $x$ a cochain of degree $p + 1$ and let
$1 \leq i_0 < \ldots < i_p \leq r$.

\medskip\noindent
Case I: $i \in \{i_0, \ldots, i_p\}$. Say $i = i_t$. Then we have
$h(d(x))_{i_0 \ldots i_p} = 0$. On the other hand we have
$$
d(h(x))_{i_0 \ldots i_p} =
\sum (-1)^j h(x)_{i_0 \ldots \hat i_j \ldots i_p} =
(-1)^t h(x)_{i_0 \ldots \hat i \ldots i_p} =
(-1)^t (-1)^{t - 1} x_{i_0 \ldots i_p}
$$
Thus $(dh + hd)(x)_{i_0 \ldots i_p} = -x_{i_0 \ldots i_p}$ as desired.

\medskip\noindent
Case II: $i \not \in \{i_0, \ldots, i_p\}$. Let $j$ be such that
$i_j < i < i_{j + 1}$. Then we see that
\begin{align*}
h(d(x))_{i_0 \ldots i_p}
& =
(-1)^j d(x)_{i_0 \ldots i_j i i_{j + 1} \ldots i_p} \\
& =
\sum\nolimits_{j' \leq j} (-1)^{j + j'}
x_{i_0 \ldots \hat i_{j'} \ldots i_j i i_{j + 1} \ldots i_p} -
x_{i_0 \ldots i_p} \\
&
+ \sum\nolimits_{j' > j} (-1)^{j + j' + 1}
x_{i_0 \ldots i_j i i_{j + 1} \ldots \hat i_{j'} \ldots i_p}
\end{align*}
On the other hand we have
\begin{align*}
d(h(x))_{i_0 \ldots i_p}
& =
\sum\nolimits_{j'} (-1)^{j'} h(x)_{i_0 \ldots \hat i_{j'} \ldots i_p} \\
& =
\sum\nolimits_{j' \leq j} (-1)^{j' + j - 1}
x_{i_0 \ldots \hat i_{j'} \ldots i_j i i_{j + 1} \ldots i_p} \\
& +
\sum\nolimits_{j' > j} (-1)^{j' + j}
x_{i_0 \ldots i_j i i_{j + 1} \ldots \hat i_{j'} \ldots i_p}
\end{align*}
Adding these up we obtain
$(dh + hd)(x)_{i_0 \ldots i_p} = - x_{i_0 \ldots i_p}$
as desired.
\end{proof}

\begin{lemma}
\label{lemma-extended-alternating-torsion}
Let $R$ be a ring. Let $f_1, \ldots, f_r \in R$. Let $M$ be an $R$-module.
Let $H^q$ be the $q$th cohomology module of the extended alternation
{\v C}ech complex of $M$. Then
\begin{enumerate}
\item $H^q = 0$ if $q \not \in [0, r]$,
\item for $x \in H^i$ there exists an $n \geq 1$ such that $f_i^n x = 0$
for $i = 1, \ldots, r$,
\item the support of $H^q$ is contained in $V(f_1, \ldots, f_r)$,
\item if there is an $f \in (f_1, \ldots, f_r)$ which acts invertibly
on $M$, then $H^q = 0$.
\end{enumerate}
\end{lemma}

\begin{proof}
Part (1) follows from the fact that the extended alternating {\v C}ech complex
is zero in degrees $< 0$ and $> r$. To prove (2) it suffices to show that
for each $i$ there exists an $n \geq 1$ such that $f_i^n x = 0$. To see
this it suffices to show that $(H^q)_{f_i} = 0$. Since localization is
exact, $(H^q)_{f_i}$ is the $q$th cohomology module of the localization
of the extended alternating complex of $M$ at $f_i$. By
Lemma \ref{lemma-extended-alternating-base-change}
this localization is the extended alternating {\v C}ech complex
of $M_{f_i}$ over $R_{f_i}$ with respect to the images of
$f_1, \ldots, f_r$ in $R_{f_i}$. Thus
we reduce to showing that $H^q$ is zero if $f_i$ is invertible, which
follows from Lemma \ref{lemma-extended-alternating-homotopy-zero}.
Part (3) follows from the observation that $(H^q)_{f_i} = 0$
for all $i$ that we just proved. To see part (4) note that in this
case $f$ acts invertibly on $H^q$ and $H^q$ is supported on $V(f)$
by (3). This forces $H^q$ to be zero (small detail omitted).
\end{proof}

\begin{lemma}
\label{lemma-extended-alternating-Cech-is-colimit-koszul}
Let $R$ be a ring. Let $f_1, \ldots, f_r \in R$. The extended alternating
{\v C}ech complex
$$
R \to \bigoplus\nolimits_{i_0} R_{f_{i_0}} \to
\bigoplus\nolimits_{i_0 < i_1} R_{f_{i_0}f_{i_1}} \to
\ldots \to R_{f_1\ldots f_r}
$$
is a colimit of the Koszul complexes $K(R, f_1^n, \ldots, f_r^n)$; see
proof for a precise statement.
\end{lemma}

\begin{proof}
We urge the reader to prove this for themselves.
Denote $K(R, f_1^n, \ldots, f_r^n)$ the Koszul complex of
Definition \ref{definition-koszul-complex} viewed as a cochain complex
sitting in degrees $0, \ldots, r$. Thus we have
$$
K(R, f_1^n, \ldots, f_r^n) :
0 \to \wedge^r(R^{\oplus r}) \to
\wedge^{r - 1}(R^{\oplus r}) \to \ldots \to
R^{\oplus r} \to R \to 0
$$
with the term $\wedge^r(R^{\oplus r})$ sitting in degree $0$.
Let $e^n_1, \ldots, e^n_r$ be the standard basis of $R^{\oplus r}$.
Then the elements $e^n_{j_1} \wedge \ldots \wedge e^n_{j_{r - p}}$ for
$1 \leq j_1 < \ldots < j_{r - p} \leq r$ form a basis for the term in
degree $p$ of the Koszul complex. Further, observe that
$$
d(e^n_{j_1} \wedge \ldots \wedge e^n_{j_{r - p}}) =
\sum (-1)^{a + 1} f_{j_a}^n e^n_{j_1} \wedge \ldots \wedge \hat e^n_{j_a} \wedge
\ldots \wedge e^n_{j_{r - p}}
$$
by our construction of the Koszul complex in Section \ref{section-koszul}.
The transition maps of our system
$$
K(R, f_1^n, \ldots, f_r^n) \to K(R, f_1^{n + 1}, \ldots, f_r^{n + 1})
$$
are given by the rule
$$
e^n_{j_1} \wedge \ldots \wedge e^n_{j_{r - p}}
\longmapsto
f_{i_0} \ldots f_{i_{p - 1}}
e^{n + 1}_{j_1} \wedge \ldots \wedge e^{n + 1}_{j_{r - p}}
$$
where the indices $1 \leq i_0 < \ldots < i_{p - 1} \leq r$ are such that
$\{1, \ldots r\} =
\{i_0, \ldots, i_{p - 1}\} \amalg \{j_1, \ldots, j_{r - p}\}$.
We omit the short computation that shows this is compatible with differentials.
Observe that the transition maps are always $1$ in degree $0$
and equal to $f_1 \ldots f_r$ in degree $r$.

\medskip\noindent
Denote $K^p(R, f_1^n, \ldots, f_r^n)$ the term of degree $p$ in the Koszul
complex. Observe that for any $f \in R$ we have
$$
R_f = \colim (R \xrightarrow{f} R \xrightarrow{f} R \to \ldots )
$$
Hence we see that in degree $p$ we obtain
$$
\colim K^p(R, f_1^n, \ldots f_r^n) =
\bigoplus\nolimits_{1 \leq i_0 < \ldots < i_{p - 1} \leq r}
R_{f_{i_0} \ldots f_{i_{p - 1}}}
$$
Here the element $e^n_{j_1} \wedge \ldots \wedge e^n_{j_{r - p}}$
of the Koszul complex above maps in the colimit
to the element $(f_{i_0} \ldots f_{i_{p - 1}})^{-n}$ in the
summand $R_{f_{i_0} \ldots f_{i_{p - 1}}}$
where the indices are chosen such that
$\{1, \ldots r\} = \{i_0, \ldots, i_{p - 1}\}
\amalg \{j_1, \ldots, j_{r - p}\}$.
Thus the differential on this complex is given by
$$
d(1\text{ in }R_{f_{i_0} \ldots f_{i_{p - 1}}}) =
\sum\nolimits_{i \not \in \{i_0, \ldots, i_{p - 1}\}}
(-1)^{i - t}\text{ in }
R_{f_{i_0} \ldots f_{i_t} f_i f_{i_{t + 1}} \ldots f_{i_{p - 1}}}
$$
Thus if we consider the map of complexes given in degree $p$
by the map
$$
\bigoplus\nolimits_{1 \leq i_0 < \ldots < i_{p - 1} \leq r}
R_{f_{i_0} \ldots f_{i_{p - 1}}}
\longrightarrow
\bigoplus\nolimits_{1 \leq i_0 < \ldots < i_{p - 1} \leq r}
R_{f_{i_0} \ldots f_{i_{p - 1}}}
$$
determined by the rule
$$
1\text{ in }R_{f_{i_0} \ldots f_{i_{p - 1}}}
\longmapsto
(-1)^{i_0 + \ldots + i_{p - 1} + p}\text{ in }R_{f_{i_0} \ldots f_{i_{p - 1}}}
$$
then we get an isomorphism of complexes from
$\colim K(R, f_1^n, \ldots, f_r^n)$ to the
extended alternating {\v C}ech complex defined in this section.
We omit the verification that the signs work out.
\end{proof}








